% CV (italiano) — Josua Aleksi Kiviranta
% Stile modellato su cv-elias-ylonen.pdf: Computer Modern classico, sezioni
% con filetto, voci in grassetto "Titolo --- Azienda" con date allineate a
% destra, paragrafi in prosa invece di elenchi puntati.
\documentclass[10pt,a4paper]{article}

\usepackage[a4paper,left=1.7cm,right=1.7cm,top=1.4cm,bottom=1.5cm]{geometry}
\usepackage[T1]{fontenc}
\usepackage[utf8]{inputenc}
\usepackage[italian]{babel}
\usepackage{lmodern}
\usepackage{microtype}
\usepackage{needspace}
\usepackage[hidelinks]{hyperref}

\pagestyle{empty}
\setlength{\parindent}{0pt}
\setlength{\parskip}{3pt}

% Titolo di sezione con filetto sottostante, come nel CV di riferimento.
\newcommand{\cvsection}[1]{%
  \vspace{9pt}%
  {\textbf{#1}}\par
  \vspace{1.5pt}\hrule\vspace{5pt}}

% Intestazione di voce: parte sinistra in grassetto, date/luogo a destra.
\newcommand{\entry}[2]{%
  \needspace{3\baselineskip}%
  \vspace{4.5pt}%
  \textbf{#1}\hfill #2\par\nopagebreak\vspace{0.5pt}}

% Riga competenze: etichetta in grassetto + descrizione.
\newcommand{\skillrow}[2]{%
  \textbf{#1} & #2 \\[1.5pt]}

\begin{document}

\begin{center}
  {\LARGE Josua Aleksi Kiviranta}\\[5pt]
  Firenze, Italia \ \textperiodcentered{} \ +358\,41\,524\,2441 \ \textperiodcentered{} \ \href{mailto:josua.kiviranta@gmail.com}{josua.kiviranta@gmail.com}
\end{center}

\cvsection{Profilo}
Ingegnere del software e imprenditore finlandese con base a Firenze. Aiuto le aziende a mettere l'intelligenza artificiale al servizio del lavoro quotidiano: automatizzare le attività ripetitive, ridurre tempi e costi dei processi e sfruttare meglio i dati. In Finlandia ho guidato team di progetto e consegnato soluzioni AI per oltre dieci grandi aziende, occupandomi sia della parte tecnica sia del rapporto con il cliente. Oggi seguo i miei clienti attraverso la mia attività di consulenza, Sassosa Consulting.

\cvsection{Esperienza}

\entry{Fondatore --- Sassosa Consulting}{2026 -- oggi \textperiodcentered{} Firenze, Italia}
Sviluppo applicazioni software e supporto le aziende nella trasformazione AI e nell'automazione dei flussi di lavoro. Tra i clienti attuali figurano aziende forestali e politici finlandesi.

\entry{AI Engineer / Project Manager --- Innofactor Oyj}{2024 -- 2025 \textperiodcentered{} Espoo, Finlandia}
Ho guidato un team di quattro persone nella consegna di soluzioni AI per oltre 10 clienti enterprise, portando i progetti dalla proposta alla produzione. Ambito tecnico: pipeline agentiche su Azure, GPT Agents ottimizzati per compiti specifici con retrieval e orchestrazione, framework di valutazione strutturati, ottimizzazione della latenza e pipeline CI/CD.

Ho contribuito alle vendite e gestito consegne di progetti end-to-end, traducendo i requisiti di business in soluzioni tecniche scalabili.

\entry{Full Stack Engineer specializzato in AI --- Innofactor Oyj}{2023 -- 2024 \textperiodcentered{} Espoo, Finlandia}
Ho diretto lo sviluppo interno di applicazioni di AI generativa e realizzato rapidamente prototipi funzionanti; ho sviluppato e lanciato \textbf{Innofactor GPT Agents} nello stesso evento nordico in cui Microsoft ha lanciato Copilot. Ho fornito consulenza sulla strategia aziendale di AI generativa, formato i colleghi sulle tecniche di AI generativa, tenuto interventi pubblici e supportato le iniziative di vendita.

\entry{Fondatore --- Pssst!}{2022 -- 2023 \textperiodcentered{} Espoo, Finlandia}
Ho costituito e guidato un team di tre persone; ho progettato e realizzato una piattaforma web per il recruiting basato sul passaparola. Ho gestito lo sviluppo del prodotto e il marketing e acquisito la base clienti iniziale.

\entry{Sviluppatore software --- Variantum Oy}{2020 \textperiodcentered{} Espoo, Finlandia}
Ho realizzato l'interfaccia utente per il sistema enterprise di gestione dei dati di prodotto di Variantum (VariPDM).

\entry{Project Manager --- Sterly Oy}{2020 \textperiodcentered{} Helsinki, Finlandia}
Ho guidato un team di 13 persone in un progetto di finanziamento ad alta urgenza in ambito governativo durante il Covid-19; ho coordinato il supporto a oltre 100 aziende finlandesi sotto forte pressione temporale.

\entry{Sottotenente --- Forze di difesa finlandesi}{2015 \textperiodcentered{} Helsinki, Finlandia}
Ho comandato un gruppo di polizia militare di 35 persone; ho sviluppato capacità di leadership e di decisione in situazioni ad alto rischio.

\cvsection{Progetti}

\entry{Verba --- \href{https://verba.world}{verba.world}}{2025}
App per l'apprendimento delle lingue basata su LLM che rileva e traduce automaticamente tutte le parole pronunciate nei video di YouTube. (React, Next.js, TypeScript, Python, Google Cloud, Azure OpenAI)

\entry{Social Synapse --- \href{https://josuakiviranta.com/socialSynapse}{socialsynapse.ai} (archiviato)}{2025}
App per aiutare le persone a fare nuove conoscenze, abbinando gli utenti in base ai loro pensieri e consumi mediatici tramite vettori semantici. (React, Node.js, TypeScript, Python, Azure Cloud, Azure OpenAI, MCP)

\entry{Ricerca di vulnerabilità informatiche --- media finlandesi}{2019}
Ho individuato e divulgato una vulnerabilità di manipolazione delle classifiche che interessava i principali quotidiani finlandesi. Ho costruito un miner online (Python) per rilevare ed esporre le manipolazioni delle classifiche e uno script proof-of-concept di manipolazione (Node.js), comunicando poi i risultati ai quotidiani interessati.

\entry{Soteduunit.fi --- \href{https://josuakiviranta.com/soteduunit}{soteduunit.fi}}{2021}
Piattaforma di annunci di lavoro basata su mappa; ho progettato e realizzato una demo per le organizzazioni finlandesi della sanità e del welfare. (React, Next.js, TypeScript)

\entry{Mining farm di criptovalute}{2017}
Ho costruito e gestito una mining farm di medie dimensioni (oltre 800\,MH/s) per finanziare i miei studi universitari.

\cvsection{Competenze}

\begin{tabular}{@{}p{2.5cm}@{\hspace{0.4cm}}p{\dimexpr\linewidth-2.9cm\relax}@{}}
\skillrow{In sintesi}{Automazione con AI, sviluppo software su misura, gestione di progetti e team}
\skillrow{AI \& LLM}{Progettazione di pipeline agentiche, RAG, context engineering, valutazione dei flussi di lavoro AI e analisi degli errori, MCP, database vettoriali, tracing, n8n, ottimizzazione della latenza}
\skillrow{Cloud}{Azure, Google Cloud, Docker, Azd, IaC, Bicep}
\skillrow{Backend}{Python, TypeScript, JavaScript, Node.js, REST API}
\skillrow{Frontend}{React, Next.js, Shadcn}
\skillrow{Dati e DB}{Pandas, NumPy, PostgreSQL + estensione vettoriale, Firebase, Supabase, MongoDB}
\skillrow{Ingegneria}{Git, GitHub Workflows (CI/CD), controllo di versione, progettazione di API}
\skillrow{Altro}{Gestione di progetti, sviluppo prodotto, collaborazione in team interfunzionali}
\end{tabular}

\cvsection{Istruzione}

\entry{Aalto University --- Laurea magistrale (M.Sc.) in Industrial Engineering \& Management}{2022 -- 2023}
Indirizzo principale: Strategia \textperiodcentered{} Indirizzo secondario: Ingegneria del software

\entry{Aalto University --- Laurea triennale (B.Sc.) in Informatica}{2017 -- 2022}
Indirizzo principale: Informatica \textperiodcentered{} Indirizzo secondario: Industrial Engineering and Management\\
Tesi: \emph{Open and linked parliamentary data and its utilisation}\\
Media: 4,02/5

\cvsection{Riconoscimenti e lingue}

\entry{Premio \emph{Best Team Player of the Year} --- Innofactor Oyj}{2024}
Riconoscimento per l'eccezionale spirito di squadra, le soluzioni innovative, la competenza tecnica e l'efficace comunicazione nelle tecnologie AI.

\entry{Lingue}{}
Finlandese (madrelingua) \textperiodcentered{} Inglese (fluente) \textperiodcentered{} Italiano (livello elementare, in miglioramento)

\cvsection{Referenze}

\vspace{2pt}
\begin{tabular}{@{}p{0.33\linewidth}@{}p{0.33\linewidth}@{}p{0.33\linewidth}@{}}
\textbf{Lasse Tolvanen} & \textbf{Vesa Syrjäkari} & \textbf{Klaus Oehlandt} \\
Co-fondatore / Business Architect & Ex vicepresidente & Membro del CdA e comproprietario \\
Innofactor Oyj & Innofactor Oyj & Variantum Oy \\
\href{mailto:lasse.tolvanen@innofactor.com}{lasse.tolvanen@innofactor.com} & \href{mailto:vesasyr@live.com}{vesasyr@live.com} & \href{mailto:klaus@oehlandt.fi}{klaus@oehlandt.fi} \\
+358 50 336 2480 & +358 50 031 6346 & +358 40 530 9060 \\
\end{tabular}

\end{document}
